\documentclass{NSF}

\begin{document}
\sloppy

\clearpage
\setcounter{page}{1}

% ===============================
\section*{Appendix 2: Facilities}
% \marginpar{\color{red} 1.0 page}
% ===============================

% ----------------------------
\subsection{Johns Hopkins University}
% ----------------------------

\begin{itemize}

\item The Open High Performance Computing cluster currently has

\begin{itemize}
\item 48,992 cores (894 nodes);
\item a combined theoretical performance of 3.9 PFLOPs;  
\item an $R_\mathrm{max}$ of 2.8 PFLOPs ($\sim \!75$\% efficiency); 
\item a combined memory space of 223 TB;
\item four high-performance file systems, GPFS and Vast, with a total of $\sim \! 20$ PB of usable space; and
\item Mellanox Infiniband HDR200 connectivity (1:1.5 topology).

\end{itemize}

\item The computing cluster is maintained by a team of 10 full-time system administrators with many years of professional experience operating high performance computer resources.

\item We have access to assistance from $\sim \! 20$ research software engineers.

\end{itemize}

% \begin{itemize}

% \item The Advanced Research Computing at Hopkins (ARCH) high-performance computing facilities (HPC) include the Rockfish cluster with over 34,000 cores (711 nodes), 144 GPUs, Mellanox Infinidad HDR100 connectivity (1:1.5 topology), and a peak performance of over 1.5 Petaflops.

% \item The Data Science and AI Institute has over 2000 servers, nearly 90,000 CPU cores, over 1100 GPUs, and over 100 petabytes of storage.

% \end{itemize}

% ----------------------------
\subsection{Cornell University}\hfill
% ----------------------------

\begin{itemize}

\item The Cornell Center for Materials Research (CCMR), \href{https://www.ccmr.cornell.edu/facilities/}{https://www.ccmr.cornell.edu/facilities/} has equipment for electron microscopy, energy-dispersive X-ray spectroscopy, and X-ray diffraction.  

\item The Cornell Center for Advanced Computing (CAC) has Red Cloud (128 core VMs/NVIDIA GPUs), 2PB of storage, and the Empire AI H100/B200 GPU clusters via a 100 Gbps campus backbone.

\item The Cornell NanoScale Science \& Technology Facility (CNF), \href{https://www.cnf.cornell.edu/equipment}{https://www.cnf.cornell.edu/equipment}, is one the nation's leading nanofabrication facilities.

\end{itemize}



% ----------------------------
\subsection{National Laboratory of the Rockies (NLR)}
% ----------------------------

\begin{itemize}
\item NLR's capabilities in photovoltaics are summarized at \href{https://www.NLR.gov/pv/index.html}{https://www.NLR.gov/pv/index.html}.
These capabilities range from fundamental studies to applied research and development (R\&D). 
Capabilities include 

\begin{itemize}
\item theory and modeling, 
\item materials deposition,
\item device design, 
\item measurements and characterization, and
\item reliability testing and engineering. 
\end{itemize}

\item NLR's photovoltaic (PV) synthesis and characterization capabilities consist of over a dozen labs comprising approximately 20,000~ft$^2$. 

\item NLR's perovskite solar-cell research program is summarized at \href{https://www.NLR.gov/pv/perovskite-solar-cells.html}{https://www.NLR.gov/pv/perovskite-solar-cells.html}.
Materials and device construction capabilities include chemical and nanomaterial synthesis; 
SOTA device fabrication, achieving \textgreater20\% efficiency in \textgreater1 cm\textsuperscript{2} cells.
\end{itemize}
% Note: Michael replaced "chemical synthesis and nanomaterial synthesis" with "chemical and nanomaterial synthesis" and "state-of-the-art" with "SOTA."

\end{document}